\documentclass[11pt,a4paper]{article}
\usepackage[utf8]{inputenc}
\usepackage[T1]{fontenc}
\usepackage[ngerman]{babel}
\usepackage[margin=0.75in]{geometry}
\usepackage{enumitem}
\usepackage{titlesec}
\usepackage{xcolor}
\usepackage{hyperref}
\definecolor{primarycolor}{RGB}{33,37,41}
\definecolor{secondarycolor}{RGB}{0,102,204}
\definecolor{accentcolor}{RGB}{108,117,125}
\hypersetup{colorlinks=true,linkcolor=secondarycolor,urlcolor=secondarycolor,citecolor=secondarycolor}
\titleformat{\section}{\large\bfseries\color{primarycolor}}{}{0em}{}
\titlespacing*{\section}{0pt}{0.8em}{0.35em}
\newcommand{\cvsection}[1]{\section*{#1}}
\setlist[itemize]{leftmargin=*,itemsep=0.05em,topsep=0.2em,parsep=0pt}
\setlength{\parindent}{0pt}
\setlength{\parskip}{0pt}
\pagestyle{empty}
\begin{document}
{\LARGE\bfseries\color{primarycolor} KATARZYNA WIŚNIEWSKA}\\[0.2em]
{\large\color{accentcolor} Pflegehelferin}\\[0.3em]
Habsburgerstraße 41, 79104 Freiburg im Breisgau\\
k.wisniewska@example.com \textbar{} +49 761 555 0412

\cvsection{PROFIL}
Pflegehelferin mit Erfahrung in der stationären Altenpflege seit 2020 und vorheriger Berufserfahrung als Krankenschwester in Polen. Sicher in der Grundpflege, qualifiziert für Behandlungspflege nach Leistungsgruppe 1 (LG1) und erfahren in der digitalen Pflegedokumentation mit Vivendi PD. Ruhig, zuverlässig und geduldig im Umgang mit Menschen mit Demenz. Deutschkenntnisse auf Niveau B2.

\cvsection{BERUFSERFAHRUNG}
\textbf{Pflegehelferin \textbar{} Seniorenzentrum Am Lindenhof, Freiburg}\\
04/2020 – heute
\begin{itemize}
\item Grundpflege für 10 bis 12 Bewohnerinnen und Bewohner pro Schicht: Körperpflege, Mobilisation, Lagerung und Hilfe beim Essen.
\item Behandlungspflege nach LG1: Blutdruck- und Blutzuckermessung, Kompressionsstrümpfe und Medikamente nach Plan reichen.
\item Digitale Pflegedokumentation in Vivendi PD.
\item Einsatz in Früh-, Spät- und Wochenenddiensten.
\end{itemize}
\smallskip
\textbf{Aushilfe Verkauf \textbar{} Bäckerei Schwarzwaldkorn, Freiburg}\\
09/2019 – 03/2020
\begin{itemize}
\item Verkauf, Kasse und Reinigung.
\end{itemize}
\smallskip
\textbf{Krankenschwester (Pielęgniarka) \textbar{} Wojewódzki Szpital Specjalistyczny, Opole, Polen}\\
10/2014 – 06/2019
\begin{itemize}
\item Tätigkeit auf einer Station für Innere Medizin mit 32 Betten.
\item Grund- und Behandlungspflege, Infusionen, Wundversorgung und Vorbereitung von Visiten.
\end{itemize}

\cvsection{AUSBILDUNG UND ANERKENNUNG}
\textbf{Bachelor Pielęgniarstwo (Krankenpflege)}\\
Medizinische Universität Opole \textbar{} 2014\\
Anerkennungsverfahren als Pflegefachfrau beim Regierungspräsidium Stuttgart: läuft.

\cvsection{WEITERBILDUNGEN}
\begin{itemize}
\item Behandlungspflege LG1 (Basiskurs), 2021
\item Expertenstandard Sturzprophylaxe (interne Schulung), 2023
\item Erste-Hilfe-Kurs, 2024
\end{itemize}

\cvsection{SPRACHEN}
Polnisch: Muttersprache\\
Deutsch: B2, telc Deutsch B1-B2 Pflege\\
Englisch: Grundkenntnisse
\end{document}